\section{Sharper recovery when few clusters are unanchored}
\label{ref:section}

The recovery construction admits a sharper analysis when its proxy error is
small.  Its polynomial running time still requires logarithmically many bad
clusters.  Throughout this section distances have the positive integral,
polynomially bounded representation defined in Section~\ref{pre:section}, and
$N=|D|+|F|$.

Theorem~\ref{ref:recovery}, stated in the introduction, is proved here.
The approximation factor has its antecedent in the FPT algorithm of
\cite[Section~2.3]{CGKLL2019}; the task here is to obtain it under the
explicit proxy promise with single-exponential dependence on the number
of bad clusters.  The proof uses the same sampling experiment and
completion algorithm as Sections~\ref{sam:section}--\ref{comp:section},
with a two-stage guide and a sharper accounting of client charges.


We use the notation $P_i=\sum_{p\in C_i}d_p^*$, and put
$\bar d_i=P_i/n_i$ when $n_i>0$.  Empty clusters have mean zero and all the
near sets below are empty.  Fix
\begin{equation}\label{ref:parameters}
 a_1=1+\mu,\qquad a_2=\frac32,\qquad
 q=1-\frac1{a_1}=\frac{\mu}{1+\mu},\qquad
 \kappa=\mu^3,
 \qquad \frac1e<\theta\le\frac1e+\mu,
\end{equation}
where $\theta$ is rational.  For $z\in\{1,2\}$ define
\[
 N_i^z=\{p\in C_i:d_p^*\le a_z\bar d_i\}.
\]
All these parameters are fixed before the input is processed.  The choice of
$\theta$ is available in Lemma~\ref{comp:partition-max}; the resulting
completion guarantee is Theorem~\ref{comp:completion}.

\subsection{A guide with two stopping rules}

Assume for now that $m>0$, and let
\[
 S_0=S\setminus\{\phi_i:i\ \mathrm{good}\},\qquad |S_0|=m.
\]
As in the sampling construction, a state consists of a set $Q\subseteq S_0$
of deleted original centers and balls indexed by a set $A$ of distinct bad
centers.  The ball for $i\in A$ has a leader $w_i\in C_i$, radius $R_i$, and
leaf dummy $\delta_i$ attached to $w_i$ with edge length $R_i$.  Write
$r_i=d(w_i,i)=d(w_i,O)$ and
\[
 H=(S\setminus Q)\cup\{\delta_i:i\in A\}\cup\{\delta_\infty\},
 \qquad B=\{i:i\ \mathrm{bad}\}\setminus A.
\]
Here $\delta_\infty$ is the permanent diameter dummy of the sampling
construction.  The computable scale and the correct radius guarantee are
\begin{equation}\label{ref:radius}
 \begin{split}
 r^0(w)&=\min\{z\ge0:
       z|\{p\in D:d(w,p)\le z\}|\ge P/m\},\\
 r_i&\le R_i\le(1+\kappa)r_i+\kappa r^0(w_i).
 \end{split}
\end{equation}
Every original-center assignment uses one fixed tie order; new dummies are
appended to that order.  Denote by $J_s(H)$ the clients assigned to
$s\in S\setminus Q$.  These assignments are defined by the state, whereas
the identities of $S_0,A,B$ are used only by the guide.

\begin{lemma}[Two-stage endpoint]\label{ref:two-stage}
The blind sampling procedure, run for at most $2m$ steps and retaining all
prefixes, has probability at least $\exp(-O_\mu(m))$ of producing a prefix
with the following properties.

There is an earlier prefix $H^0$, with covered set $A^1$ and remaining bad set
$B^0$, at which the deleted set is the final set $Q$.  Put
\[
 R^*=S_0\setminus Q,\qquad
 J^0=\bigcup_{s\in R^*}J_s(H^0),\qquad
 T_i=\sum_{p\in N_i^1}d(p,H^0)\quad(i\in B^0).
\]
Then
\begin{equation}\label{ref:first-stop}
 \sum_{i\in B^0}T_i\le\kappa P,
 \qquad \sum_{p\in J^0}d(p,H^0)\le\kappa P.
\end{equation}
For each nonempty cluster $i\in B^0$, define
$f_i=|C_i\cap J^0|/n_i$.  All balls added after $H^0$ correspond to a set
$A^2\subseteq\{i\in B^0:n_i>0,\ f_i<1/3\}$, and their leaders lie in
$N_i^2$; the leaders for $A^1$ lie in $N_i^1$.  At the final prefix, set
$A=A^1\cup A^2$, $B=B^0\setminus A^2$, and
$J=\bigcup_{s\in R^*}J_s(H)$.  We have
\begin{equation}\label{ref:second-stop}
 J\subseteq J^0,\qquad
 \sum_{p\in J}d(p,H)\le\kappa P,\qquad
 \sum_{\substack{i\in B:n_i>0\\ f_i<1/3}}
       \sum_{p\in N_i^2}d(p,H)\le\kappa P.
\end{equation}
Every ball satisfies Equation~\eqref{ref:radius}, and the collection satisfies
\begin{equation}\label{ref:radius-sum}
 \sum_{i\in A}n_i r^0(w_i)\le\frac{11}{2}P.
\end{equation}
The earlier prefix and the conditions defining it need not be recognized by
the algorithm.
\end{lemma}

\begin{proof}
The proof specifies a guide inside the same blind sampling procedure.
For every prescribed mode and eligible client set, the guide draws from
the actual law conditioned on that event and uses the correct radius-menu
entry specified in Lemma~\ref{sam:radius-budget}.
At a state of its first stage, write
\[
 T(H)=\sum_{i\in B}\sum_{p\in N_i^1}d(p,H),\qquad
 I(H)=\sum_{s\in S_0\setminus Q}\sum_{p\in J_s(H)}d(p,H).
\]
If $T(H)>\kappa P$, the guide chooses a ball step whose sampled client belongs
to $N_i^1$ for some $i\in B$.  If $T(H)\le\kappa P$ but
$I(H)>\kappa P$, it chooses a removal step whose sampled client is assigned
to a member of $S_0\setminus Q$.  It stops this stage as soon as both
inequalities hold in the other direction.

Here are the probability bounds needed to invoke the general guide lemma.
Lemma~\ref{sam:near-clients} gives $|N_i^1|\ge qn_i$ for every nonempty
uncovered cluster, and hence
\begin{equation}\label{ref:near-bound}
 d(i,H)\le a_1\bar d_i+
       \frac{\sum_{p\in N_i^1}d(p,H)}{qn_i}.
\end{equation}
Indeed, take a member of $N_i^1$ of minimum distance to $H$ and use the
triangle inequality.  Summing the resulting bound for all clients of an
uncovered cluster gives
\[
 \sum_{p\in C_i}d(p,H)
 \le(1+a_1)P_i+\frac1q\sum_{p\in N_i^1}d(p,H).
\]
Good clients have their proxies available.  A client of a covered cluster can
use its dummy at distance at most
$d_p^*+(2+\kappa)a_2\bar d_i+\kappa r^0(w_i)$.
Since every leader used in either stage has $r_i\le a_2\bar d_i$,
Lemma~\ref{sam:computable-radius} bounds the sum of the last radius terms by
$(11/2)\kappa P$.  Thus throughout the first stage the total sampling weight
$W(H)=\sum_{p\in D}d(p,H)$ satisfies the convenient bound
\begin{equation}\label{ref:weight}
 W(H)\le 6P+T(H)/q.
\end{equation}
For example, the coefficient of cluster cost before the additive proxy and
radius errors is at most $1+(2+\kappa)a_2$, and
$1+(2+\kappa)a_2+\mu+(11/2)\kappa<6$ for
$0<\mu\le1/4$.

The procedure samples proportionally to this weight and chooses its mode
with probability $1/2$.  When $T(H)>\kappa P$, the joint event prescribed by
the guide has probability at least
\[
 \frac{T(H)}{2(6P+T(H)/q)}
 \ge \rho_\mu:=\frac{\kappa}{2(6+\kappa/q)}>0.
\]
When only the removal condition fails, its joint event has the same lower
bound by Equation~\eqref{ref:weight}.  A zero total weight already satisfies both
stopping conditions.  At the first stop call the state $H^0$ and use the
notation of the statement; this proves Equation~\eqref{ref:first-stop}.

In the second stage the guide deletes no further original centers.  The
eligible centers are fixed using $H^0$: they are precisely the nonempty
$i\in B^0$ with $f_i<1/3$.  While the final sum in
Equation~\eqref{ref:second-stop} exceeds $\kappa P$, the guide adds a ball whose
leader is in $N_i^2$ for an eligible, still uncovered center $i$.
All reference distances can only decrease after $H^0$, so
\[
 W(H)\le W(H^0)\le(6+\kappa/q)P.
\]
The prescribed joint event therefore again has probability at least
$\rho_\mu$.  Its termination gives the last inequality in
Equation~\eqref{ref:second-stop}.

The reference set only grows during this stage, and the order between old
choices is fixed.  A client that was assigned outside the fixed set $R^*$
cannot become assigned to it upon adding choices.  Consequently
$J\subseteq J^0$, and monotonicity of distances gives the internal-cost
bound in Equation~\eqref{ref:second-stop}.

Across both stages each step either deletes a fresh member of $S_0$ or
covers a fresh bad center.  Thus there are at most $2m$ steps in total.
All leaves persist across the transition.  Lemma~\ref{sam:radius-budget}
therefore applies once to the entire trajectory: its disjoint-ball
amortization bounds the sum of costly radius-menu indices across both
stages.  Correct radii satisfy Equation~\eqref{ref:radius}, and
Lemma~\ref{sam:computable-radius} gives Equation~\eqref{ref:radius-sum} because
$a_1\le a_2=3/2$.  The joint progress probabilities just proved and the
global radius budget meet the hypotheses of Lemma~\ref{sam:probability}.
That lemma gives the claimed $\exp(-O_\mu(m))$ probability.
In particular, its guide may switch stages using the unobserved state
$H^0$: the actual procedure continues its ordinary sampling and mode
choices and retains every prefix.
\end{proof}

\subsection{The completion charges}

Fix a successful endpoint from Lemma~\ref{ref:two-stage}.  Apply
Theorem~\ref{comp:completion} with parameter $\theta$ and the surrogate
$\Psi$ defined there.  Its comparison choice in the copied ball parts is
$X^*=\{(i,i):i\in A\}$, with distances evaluated at the underlying
facilities.  Conditional on its successful colors and labels, it
returns a solution of cardinality at most $h$, and its expected cost is at
most
\begin{equation}\label{ref:completion-bound}
 2I+\sum_{p\in J}c_p+\sum_{p\notin J}b_p,
 \qquad I=\sum_{p\in J}d(p,H),
\end{equation}
where, for $p\in C_i$,
\[
 c_p=
 \begin{cases}
   d(p,\phi_i),&i\ \mathrm{good},\\
   d_p^*+\theta(r_i+R_i),&i\in A,\\
   d_p^*,&i\in B,
 \end{cases}
 \qquad
 b_p=\theta d(p,H)+(1-\theta)
       \min\{d(p,H),d(p,X^*)\}.
\]
As usual a distance to the empty set is $+\infty$.  Keeping the exact
expression for $b_p$ permits different estimates on the two groups of
covered clusters.

\begin{lemma}[Refined charge bound]\label{ref:charges}
For the endpoint above, the cost bound in
Equation~\eqref{ref:completion-bound} is at most
\begin{equation}\label{ref:explicit-bound}
 \left(1+2\theta a_1+\mu+11\kappa+\frac\kappa q\right)P
 \le (1+2/e+6\mu)P.
\end{equation}
\end{lemma}

\begin{proof}
For each cluster let $C_i^{\mathrm{chg}}$ denote the sum of its $c_p$ terms
on $J$ and its $b_p$ terms outside $J$.  We account for these charges first,
and then for $2I$.

\paragraph{Good clusters.}
Their proxies remain in $H$, so $b_p\le d(p,H)\le d(p,\phi_i)$.
Together with the definition of $c_p$, their total charge is at most
$P_g+\mu P$.

\paragraph{Clusters covered in the first stage.}
For $i\in A^1$, the comparison choice $X^*$ offers center $i$, and its dummy gives
$d(p,H)\le d_p^*+r_i+R_i$.  Thus both types of client charge are at most
$d_p^*+\theta(r_i+R_i)$.  Since $r_i\le a_1\bar d_i$,
\begin{equation}\label{ref:first-covered}
 C_i^{\mathrm{chg}}
 \le (1+2\theta a_1)P_i
       +\theta\kappa a_1P_i
       +\theta\kappa n_i r^0(w_i).
\end{equation}

\paragraph{Clusters covered in the second stage.}
For a nonempty $i\in A^2$, put
$x_i=|C_i\cap J|/n_i\le f_i<1/3$.  On $J$ use the ball charge with
$r_i\le a_2\bar d_i$.  Outside $J$, use the center $i$ offered by $X^*$ and the
first-stage estimate in Equation~\eqref{ref:near-bound} at $H^0$:
\[
 b_p\le d_p^*+\theta\left(a_1\bar d_i+\frac{T_i}{qn_i}\right).
\]
It follows that
\begin{align}
 C_i^{\mathrm{chg}}
 &\le P_i+\theta\bigl(2a_2x_i+(1-x_i)a_1\bigr)P_i
       +\theta\kappa a_2x_iP_i
       +\theta\kappa x_in_i r^0(w_i)
       +\frac{\theta(1-x_i)T_i}{q}\notag\\
 &\le (1+2\theta)P_i
       +\theta\kappa a_2x_iP_i
       +\theta\kappa x_in_i r^0(w_i)
       +\frac{T_i}{q}.
 \label{ref:second-covered}
\end{align}
For the last inequality, $2a_2-a_1=2-\mu>0$ and
\begin{equation}\label{ref:second-coefficient}
 2a_2x_i+(1-x_i)a_1
 \le 3f_i+(1-f_i)(1+\mu)
 \le \frac53+\frac{2\mu}{3}\le2.
\end{equation}
In particular this coefficient estimate uses the fraction at the first
stop and the inclusion $J\subseteq J^0$.

\paragraph{Remaining clusters with $f_i\ge1/3$.}
Clients in $J$ have charge $d_p^*$.  Clients in $J^0\setminus J$ have
$b_p\le d(p,H)\le d(p,H^0)$.  For the at most $(2/3)n_i$ clients outside
$J^0$, use
$b_p\le d(p,H)\le d_p^*+a_1\bar d_i+T_i/(qn_i)$.
Consequently
\begin{equation}\label{ref:uncovered-large}
 C_i^{\mathrm{chg}}
 \le \left(1+\frac23a_1\right)P_i+\frac{T_i}{q}+E_i^+,
 \qquad
 E_i^+=\sum_{p\in C_i\cap(J^0\setminus J)}d(p,H^0).
\end{equation}
The sum of $E_i^+$ over this group of clusters is at most $\kappa P$ by
Equation~\eqref{ref:first-stop}.

\paragraph{Remaining clusters with $f_i<1/3$.}
Again clients in $J$ have charge $d_p^*$.  For clients outside $J$ that lie
in $N_i^2$, retain the base bound $b_p\le d(p,H)$.  Every other client
outside $J$ has $d_p^*>a_2\bar d_i$, whence
\[
 b_p\le d_p^*+a_1\bar d_i+\frac{T_i}{qn_i}
 \le\left(1+\frac23a_1\right)d_p^*+\frac{T_i}{qn_i}.
\]
Thus
\begin{equation}\label{ref:uncovered-small}
 C_i^{\mathrm{chg}}
 \le\left(1+\frac23a_1\right)P_i+\frac{T_i}{q}+E_i^-,
 \qquad
 E_i^-=\sum_{p\in (C_i\cap N_i^2)\setminus J}d(p,H).
\end{equation}
The sum of $E_i^-$ over these remaining eligible clusters is at most
$\kappa P$ by Equation~\eqref{ref:second-stop}.  Empty clusters have zero charge and
require none of the divisions above.

\paragraph{Summing the errors once.}
The main coefficient for every bad cluster is bounded by $1+2\theta a_1$:
this follows from $a_1\ge1$ and $\theta>1/e>1/3$.
The $T_i/q$ terms in Equations~\eqref{ref:second-covered},
\eqref{ref:uncovered-large}, and~\eqref{ref:uncovered-small} concern
disjoint subsets of $B^0$.  Their total is therefore at most $\kappa P/q$.
The radius errors from Equations~\eqref{ref:first-covered} and
\eqref{ref:second-covered} together are at most
\[
 \theta\kappa\left(
   a_1\sum_{i\in A^1}P_i+a_2\sum_{i\in A^2}x_iP_i
   +\sum_{i\in A^1}n_i r^0(w_i)
   +\sum_{i\in A^2}x_in_i r^0(w_i)\right)
 \le 7\theta\kappa P\le7\kappa P.
\]
Here $a_1\le a_2=3/2$, $x_i\le1$, and Equation~\eqref{ref:radius-sum} are used
globally.  The two negligible base-cost sums contribute at most $2\kappa P$,
and $2I\le2\kappa P$. Along with the proxy error, these bounds establish
the bound by the first expression in Equation~\eqref{ref:explicit-bound}.

Finally, Equation~\eqref{ref:parameters} implies
\[
 2\theta a_1+\mu+11\kappa+\frac\kappa q
 \le \frac2e+\left(3+\frac2e\right)\mu
                +3\mu^2+12\mu^3
 \le \frac2e+6\mu.
\]
For the last step, $\mu\le1/4$ gives
$3\mu^2+12\mu^3\le(3/2)\mu$, and $3+2/e+3/2<6$.
In particular the cumulative sampling and radius errors are
$O(\kappa+\kappa/q)P=O(\mu^2)P$; the linear error comes from the proxies,
$a_1-1$, and $\theta-1/e$.
\end{proof}

\subsection{Extraction, running time, and scope}

\begin{proof}[Proof of Theorem~\ref{ref:recovery}]
The algorithm retains $S$ as a feasible candidate.  When $m=0$,
Equation~\eqref{ref:proxy-hypothesis} gives
$\operatorname{cost}(S)\le(1+\mu)P$, which suffices.
For $m>0$, first consider the branch with the correct guesses of $P$ and
$m$.  Lemma~\ref{ref:two-stage} produces a suitable prefix with probability
at least $\exp(-O_\mu(m))$.  Completion is attempted at every prefix.
Theorem~\ref{comp:completion} supplies the needed colors and labels with
probability $\exp(-O(m))$ and always respects the cardinality bound.
Conditional on these choices, Lemma~\ref{ref:charges} gives
\[
 \mathbb E[\operatorname{cost}(\widehat S)]
 \le(1+2/e+6\mu)P
 \le(1+2/e+\varepsilon/2)P.
\]
By nonnegativity of cost, Markov's inequality therefore gives conditional
probability at least
\[
 1-\frac{1+2/e+\varepsilon/2}{1+2/e+\varepsilon}
 =\frac{\varepsilon}{2(1+2/e+\varepsilon)}>0
\]
of meeting the stated target.  One full trial consequently succeeds with
probability $\exp(-O_{\mu,\varepsilon}(m))$.

If $\Delta$ is the largest distance in the scaled input, then the unknown
integer $P$ belongs to $\{1,\ldots,n\Delta\}$.  Try all these values and
all $m\le\min\{h,\lfloor L_0\log N\rfloor\}$.
The sampling menus are finite and have polynomial bit complexity by
Lemmas~\ref{sam:computable-radius} and~\ref{sam:radius-budget}; their
single-exponential probability cost has already been accounted for over
the whole trajectory.  Each completion uses
$\exp(O(m))\operatorname{poly}(N)$ work for fixed $\theta$, by
Theorem~\ref{comp:completion}.  Repeating the trials
$\exp(O_{\mu,\varepsilon}(L_0\log N))\log(1/\zeta)$ times, with the
constant enlarged if necessary, makes the failure probability at most
$\zeta$.  All enumerations, repetitions, and prefix completions therefore
take polynomial time for the fixed parameters.  Take the cheapest feasible
candidate; no successful guide, stopping state, coloring, or comparison
solution needs to be identified.
\end{proof}

The proxy accuracy in Theorem~\ref{ref:recovery} is supplied by the
parameterized anchoring theorem whenever its hypotheses hold.  Specifically,
Theorem~\ref{loc:anchoring} gives error $\mu P$ and bad count at most
$C_\mu P/b$, with the finite constant
\[
 C_\mu=10+\frac{600}{\mu}\,2^{\lceil40/\mu\rceil-1}.
\]
Thus if $P/b=O(\log N)$, its output satisfies the required logarithmic
bound, with a constant depending on $\mu$.  Declaring every comparison
cluster bad also gives the hypothesis when $h=O(\log N)$.
The low-sensitivity part of the main argument supplies no logarithmic bound
on this bad count.  Its guarantee therefore remains the fixed improvement
below $2$ proved there; Theorem~\ref{ref:recovery} sharpens the recovery
component within the scope of Equation~\eqref{ref:proxy-hypothesis}.
